```latex
\documentclass[11pt]{amsart}

\usepackage{amsmath,amssymb,amsthm}
\usepackage{hyperref}

\newtheorem{theorem}{Theorem}[section]
\newtheorem{proposition}[theorem]{Proposition}
\newtheorem{lemma}[theorem]{Lemma}
\theoremstyle{definition}
\newtheorem{definition}[theorem]{Definition}
\theoremstyle{remark}
\newtheorem{remark}[theorem]{Remark}

\DeclareMathOperator{\NS}{NS}
\DeclareMathOperator{\Pic}{Pic}
\DeclareMathOperator{\rank}{rank}
\DeclareMathOperator{\disc}{disc}
\DeclareMathOperator{\sgn}{sgn}

\newcommand{\cX}{\mathcal{X}}
\newcommand{\cL}{\mathcal{L}}
\newcommand{\D}{\Delta}
\newcommand{\Z}{\mathbb{Z}}
\newcommand{\C}{\mathbb{C}}
\newcommand{\PP}{\mathbb{P}}

\begin{document}

\title{Mirror pairs for $H$-polarized Tyurin degenerations}
\author{}
\date{}
\maketitle

\section{Statement}

Let $H$ be the hyperbolic plane, and let $E_8$ denote the negative definite
$E_8$ lattice.  We use the corrected definition of an $H$-polarized Tyurin
degeneration stated in the question.  In particular, the ample model class $A$
is required to satisfy
\[
(A|_{V_i})^2>0
\qquad (i=1,2).
\]
This condition is used in the mirror-pair criterion and is not a consequence
of the other assumptions.

\begin{theorem}\label{thm:main}
Let
\[
\pi\colon \cX\longrightarrow \D
\]
be a Tyurin degeneration of $K3$ surfaces whose central fibre is
\[
X_0=V_1\cup V_2,
\]
and suppose that it is $H$-polarized in the corrected sense of
\cite[Definition~3.20]{dfk3slpms}.  Assume that the image of $0$ in the
Alexeev--Brunyate--Engel compactification of the moduli space of
$H$-polarized $K3$ surfaces is not contained in any Noether--Lefschetz locus.

Let $Y$ be a very general
\[
M:=H\oplus E_8\oplus E_8
\]
quasipolarized $K3$ surface.  Then there exist
\begin{enumerate}
\item an elliptic fibration
\[
\pi_Y\colon Y\longrightarrow \PP^1;
\]
\item an allowable loop $\gamma\subset \PP^1$ in the sense of
\cite[Definition~4.11]{dfk3slpms},
\end{enumerate}
such that the $M$-quasipolarization is compatible with $\pi_Y$ and $\gamma$
in the sense of \cite[Definition~4.21]{dfk3slpms}, and such that
\[
\bigl(\cX/\D,\; Y/\PP^1,\gamma\bigr)
\]
is a mirror pair in the sense of \cite[Definition~5.1]{dfk3slpms}.
\end{theorem}

The proof has three parts.  First, we identify the mirror lattice.  Second,
we realize the $M$-quasipolarization by an elliptic fibration with an
allowable loop.  Finally, we verify that the hypotheses of the
Giovenzana--Thompson mirror-pair criterion are satisfied.

\section{The mirror lattice}

Let
\[
L_{K3}\simeq H^{\oplus 3}\oplus E_8^{\oplus 2}
\]
be the even unimodular $K3$ lattice.

\begin{lemma}\label{lem:M-lattice}
The lattice
\[
M=H\oplus E_8\oplus E_8
\]
is even and unimodular, and has signature $(1,17)$.
\end{lemma}

\begin{proof}
The hyperbolic plane is even, has signature $(1,1)$, and satisfies
$\disc H=-1$.  The negative $E_8$ lattice is even, unimodular, and has
signature $(0,8)$.  Therefore the displayed orthogonal direct sum is even,
has determinant $-1$, and has signature
\[
(1,1)+(0,8)+(0,8)=(1,17).
\]
Thus $M$ is unimodular.
\end{proof}

The following lattice calculation explains the origin of the mirror
polarization.

\begin{lemma}\label{lem:cusp-quotient}
Let $I\subset L_{K3}$ be a primitive totally isotropic sublattice of rank
two.  Then
\[
I^\perp/I\simeq H\oplus E_8\oplus E_8.
\]
\end{lemma}

\begin{proof}
Because $I$ is primitive and totally isotropic, the induced quadratic form on
$I^\perp/I$ is even and nondegenerate.  Since $L_{K3}$ is unimodular and $I$
is primitive, the quotient $I^\perp/I$ is also unimodular; this is the
standard discriminant calculation for primitive isotropic sublattices
\cite{nikulin}.  Its rank is
\[
22-2\rank I=18.
\]
Passing from $L_{K3}$ to $I^\perp/I$ removes one hyperbolic plane from the
signature, so
\[
\sgn(I^\perp/I)=(3,19)-(1,1)=(1,17).
\]
Indefinite even unimodular lattices are classified by their signatures
\cite{nikulin}.  Since
\[
17-1=16\equiv 0\pmod 8,
\]
the unique even unimodular lattice of signature $(1,17)$ is
\[
H\oplus E_8\oplus E_8.
\]
This proves the lemma.
\end{proof}

\begin{remark}
For a Type~III degeneration of $K3$ surfaces, the limiting mixed Hodge
structure has a primitive totally isotropic plane $I=W_0$; see
\cite{schmid,kulikov}.  Lemma~\ref{lem:cusp-quotient} identifies the
associated mirror polarization lattice as $M$.
\end{remark}

We next record the usual description of a very general lattice-polarized
$K3$ surface.

\begin{lemma}\label{lem:very-general-NS}
If $Y$ is a very general $M$-quasipolarized $K3$ surface, then the given
primitive embedding $M\hookrightarrow \NS(Y)$ is an isomorphism:
\[
\NS(Y)\simeq M.
\]
\end{lemma}

\begin{proof}
Let $\omega_Y$ be the period of $Y$.  If
$d\in H^2(Y;\Z)$ represents an integral $(1,1)$-class not lying in $M$, then
\[
d\cdot\omega_Y=0.
\]
For fixed $d\notin M$, this equation cuts out a proper Noether--Lefschetz
locus in the $M$-polarized period domain.  There are countably many such
classes.  A very general point avoids their union, and hence every integral
$(1,1)$-class lies in the image of $M$.  Thus $\NS(Y)=M$.
\end{proof}

\section{The compatible elliptic fibration}

We now use the standard elliptic criterion for $K3$ surfaces, together with
the compatible-loop construction of Giovenzana--Thompson.

\begin{proposition}\label{prop:compatible-fibration}
Let $Y$ be a very general $M$-quasipolarized $K3$ surface.  Then there exist
an elliptic fibration with section
\[
\pi_Y\colon Y\longrightarrow \PP^1
\]
and an allowable loop $\gamma\subset\PP^1$ such that the given
$M$-quasipolarization is compatible with $\pi_Y$ and $\gamma$ in the sense of
\cite[Definition~4.21]{dfk3slpms}.

Moreover, the compatible realization may be chosen so that:
\begin{enumerate}
\item the $H$-summand is generated by the fibre and section classes;
\item the two $E_8$-summands occur as the root lattices of the two
distinguished reducible fibres lying on opposite sides of $\gamma$.
\end{enumerate}
\end{proposition}

\begin{proof}
By Lemma~\ref{lem:very-general-NS}, $\NS(Y)\simeq M$.  The hyperbolic plane
contains a primitive isotropic vector $f$.  After replacing $f$ by $-f$ and
changing the chosen Weyl chamber, if necessary, $f$ lies in the closure of the
ample cone.  A primitive nef isotropic divisor on a $K3$ surface is
base-point free, and the morphism associated to $|f|$ is an elliptic
fibration; see \cite{saintdonat} and \cite[Chapter~2]{huybrechts}.  Since
$\NS(Y)$ contains the hyperbolic plane, the fibration can be chosen with a
section: this is the standard $U$-lattice criterion for elliptic $K3$
surfaces with section \cite[Chapter~11]{huybrechts}.

Let $s$ and $f$ denote the section and fibre classes.  Then
\[
s^2=-2,\qquad f^2=0,\qquad s\cdot f=1,
\]
so $\langle s,f\rangle\simeq H$.  In the unimodular lattice
\[
M=\langle s,f\rangle\oplus E_8\oplus E_8,
\]
the root frame orthogonal to $\langle s,f\rangle$ is $E_8\oplus E_8$.
Each $E_8$-factor is irreducible as a root system and therefore corresponds
to one reducible fibre whose dual graph is affine $E_8$; by the Kodaira
classification, such a fibre is of type $\mathrm{II}^*$.

The existence of an allowable loop separating these two distinguished fibres,
and the fact that the loop and fibration may be chosen compatibly with the
fixed $M$-quasipolarization, is the rank-two case of the compatible
fibration construction in \cite[\S4]{dfk3slpms}.  The loop is taken in the
complement of the discriminant locus of $\pi_Y$, and it is chosen so that
the two $E_8$-factors lie in different connected components of
$\PP^1\setminus\gamma$.  This is precisely the lattice compatibility
required by \cite[Definition~4.21]{dfk3slpms}.
\end{proof}

\begin{remark}
The remaining singular fibres, if any, do not affect the compatibility
condition.  The allowable loop is chosen after all critical values have been
fixed, and the loop criterion in \cite[\S4]{dfk3slpms} permits a generic
embedded loop separating the two distinguished $E_8$-fibres.
\end{remark}

\section{The mirror-pair criterion}

We now state the mirror correspondence input from
\cite{dfk3slpms}.  This is the point at which the corrected positivity
condition on both components of the central fibre is used.

\begin{theorem}[Giovenzana--Thompson mirror-pair criterion]
\label{thm:GT-criterion}
Let $\cX\to\D$ be a Tyurin degeneration satisfying the corrected definition
of an $H$-polarization, including
\[
(A|_{V_1})^2>0,\qquad (A|_{V_2})^2>0.
\]
Suppose that the image of $0$ in the Alexeev--Brunyate--Engel
compactification is not contained in any Noether--Lefschetz locus.

Let $Y$ be a very general
\[
H\oplus E_8\oplus E_8
\]
quasipolarized $K3$ surface, and let $\pi_Y\colon Y\to\PP^1$ and
$\gamma\subset\PP^1$ be compatible fibration data as in
Proposition~\ref{prop:compatible-fibration}.  Then
\[
\bigl(\cX/\D,\;Y/\PP^1,\gamma\bigr)
\]
is a mirror pair in the sense of \cite[Definition~5.1]{dfk3slpms}.
\end{theorem}

\begin{proof}
This is the rank-two case of the mirror-pair correspondence proved in
\cite[\S5]{dfk3slpms}.  For clarity, we indicate how the hypotheses enter.

The limiting mixed Hodge structure of the Type~III degeneration determines a
primitive totally isotropic plane $I\subset L_{K3}$.  By
Lemma~\ref{lem:cusp-quotient},
\[
I^\perp/I\simeq H\oplus E_8\oplus E_8,
\]
which is the polarization lattice on the fibration side.  The
Alexeev--Brunyate--Engel compactification identifies the corresponding
boundary point with the stable-pair data of the two components
$(V_1,D_1)$ and $(V_2,D_2)$ \cite{cmek3sspt}.  The corrected inequalities
\[
(A|_{V_i})^2>0
\]
ensure that the classes induced by $A$ satisfy the positivity condition
required on both components in the mirror-pair comparison.

Finally, the assumption that the boundary point lies in no
Noether--Lefschetz locus implies that no additional integral $(1,1)$-class
occurs in the limiting period.  Under this very-generality assumption, the
comparison theorem of \cite[\S5]{dfk3slpms} identifies the monodromy and
lattice data on the two sides of the allowable loop with the two components
of the Type~III central fibre.  These are the conditions defining a mirror
pair in \cite[Definition~5.1]{dfk3slpms}.
\end{proof}

\section{Proof of the main theorem}

We now assemble the preceding results.

\begin{proof}[Proof of Theorem~\ref{thm:main}]
By assumption, $\cX\to\D$ satisfies the corrected conditions of
\cite[Definition~3.20]{dfk3slpms}.  In particular:
\begin{enumerate}
\item the local subsystem $\cL$ with fibre $H$ is monodromy invariant and
defines an $H$-quasipolarization on every smooth fibre;
\item the effective flat nef Cartier divisor $A$ restricts to a nef and big
class in $\cL_p$ for every $p\in\D^*$;
\item the two numerical inequalities
\[
(A|_{V_1})^2>0,\qquad (A|_{V_2})^2>0
\]
hold.
\end{enumerate}
Moreover, the image of $0$ in the Alexeev--Brunyate--Engel compactification
avoids every Noether--Lefschetz locus.

Let $Y$ be the fixed very general
$H\oplus E_8\oplus E_8$-quasipolarized $K3$ surface.  By
Proposition~\ref{prop:compatible-fibration}, there exist an elliptic
fibration
\[
\pi_Y\colon Y\longrightarrow \PP^1
\]
and an allowable loop $\gamma\subset\PP^1$ with which the given
quasipolarization is compatible in the sense of
\cite[Definition~4.21]{dfk3slpms}.

The degeneration and these fibration data now satisfy all hypotheses of the
mirror-pair criterion in Theorem~\ref{thm:GT-criterion}.  Therefore
\[
\bigl(\cX/\D,\;Y/\PP^1,\gamma\bigr)
\]
is a mirror pair in the sense of \cite[Definition~5.1]{dfk3slpms}.
\end{proof}

\begin{remark}[Role of the corrected definition]
The final inequalities in condition (2) are essential in the application of
Theorem~\ref{thm:GT-criterion}.  They guarantee that the pseudo-ample class
induced by $A$ has positive self-intersection on both components $V_1$ and
$V_2$.  Thus the corrected definition, rather than the incomplete arXiv
version, is the one used in the proof.
\end{remark}

\begin{thebibliography}{ABE22}

\bibitem[ABE22]{cmek3sspt}
V.~Alexeev, A.~Brunyate, and P.~Engel,
\emph{Compactifications of moduli of elliptic $K3$ surfaces: Stable pair and
toroidal},
Geom. Topol. \textbf{26} (2022), no.~8, 3525--3588.

\bibitem[GT24]{dfk3slpms}
L.~Giovenzana and A.~Thompson,
\emph{Degenerations and fibrations of $K3$ surfaces: Lattice polarisations and
mirror symmetry},
Preprint, May 2024, arXiv:2405.12009.

\bibitem[Huy16]{huybrechts}
D.~Huybrechts,
\emph{Lectures on $K3$ surfaces},
Cambridge Studies in Advanced Mathematics \textbf{158},
Cambridge University Press, Cambridge, 2016.

\bibitem[Kul77]{kulikov}
V.~S.~Kulikov,
\emph{Degenerations of $K3$ surfaces and Enriques surfaces},
Izv. Akad. Nauk SSSR Ser. Mat. \textbf{41} (1977), no.~5, 1008--1042;
English translation, Math. USSR-Izv. \textbf{11} (1977), 957--989.

\bibitem[Nik80]{nikulin}
V.~V.~Nikulin,
\emph{Integer symmetric bilinear forms and some of their geometric
applications},
Izv. Akad. Nauk SSSR Ser. Mat. \textbf{43} (1979), 111--177;
English translation, Math. USSR-Izv. \textbf{14} (1980), 103--167.

\bibitem[SD74]{saintdonat}
B.~Saint-Donat,
\emph{Projective models of $K3$ surfaces},
Amer. J. Math. \textbf{96} (1974), 602--639.

\bibitem[Sch73]{schmid}
W.~Schmid,
\emph{Variation of Hodge structure: the singularities of the period mapping},
Invent. Math. \textbf{22} (1973), 211--319.

\end{thebibliography}

\end{document}
```